% ===== PRÉAMBULE CANONIQUE — à coller VERBATIM en tête de CHAQUE .tex =====
\documentclass[11pt,a4paper]{article}
\usepackage[utf8]{inputenc}\usepackage[T1]{fontenc}\usepackage{lmodern}
\usepackage[french]{babel}
\usepackage{amsmath,amssymb,mathtools}\usepackage{icomma}
\usepackage[version=4]{mhchem}          % \ce{...} : équations & formules
\usepackage{siunitx}\sisetup{locale=FR} % \qty{}{}, \si{}, \ang{}
\usepackage{chemfig}                    % \chemfig{...} : structures développées
\usepackage{xcolor}\usepackage{colortbl}
\usepackage{tikz}\usepackage{pgfplots}\pgfplotsset{compat=1.18}
\usepackage[most]{tcolorbox}\usepackage{enumitem}\usepackage{array,booktabs}
\usepackage[a4paper,margin=2.2cm]{geometry}
\usetikzlibrary{arrows.meta,calc,angles,quotes,positioning,patterns,backgrounds,shapes.geometric}
\definecolor{primary}{RGB}{222,49,99}\definecolor{primarydark}{RGB}{150,22,55}
\definecolor{defcol}{RGB}{0,114,178}\definecolor{methcol}{RGB}{52,92,148}
\definecolor{excol}{RGB}{0,135,90}\definecolor{attncol}{RGB}{200,40,55}
\tcbset{boxbase/.style={enhanced,breakable,boxrule=0.9pt,arc=2.5pt,
  left=3mm,right=3mm,top=2.3mm,bottom=2.3mm,fonttitle=\bfseries,coltitle=white}}
% --- boîtes TUTORIEL (noms EXACTS, ne pas renommer) ---
\newtcolorbox{cle}[1][]{boxbase,colback=defcol!6,colframe=defcol,
  title={\textsc{À retenir}\ifx\relax#1\relax\relax\else\textmd{ : \itshape #1}\fi}}
\newtcolorbox{methode}[1][]{boxbase,colback=methcol!7,colframe=methcol,sharp corners,
  title={\textsc{Méthode}\ifx\relax#1\relax\relax\else\textmd{ --- \itshape #1}\fi}}
\newtcolorbox{exemplebox}[1][]{boxbase,colback=excol!5,colframe=excol,
  title={\textsc{Exemple}\ifx\relax#1\relax\relax\else\textmd{ : \itshape #1}\fi}}
\newtcolorbox{piege}[1][]{boxbase,colback=attncol!6,colframe=attncol,
  title={\textsc{Piège}\ifx\relax#1\relax\relax\else\textmd{ : \itshape #1}\fi}}
\newtcolorbox{remarque}[1][]{boxbase,colback=black!4,colframe=black!45,
  title={\textsc{Remarque}\ifx\relax#1\relax\relax\else\textmd{ : \itshape #1}\fi}}
% --- boîtes EXERCICE / CORRIGÉ (noms EXACTS, auto-comptées) ---
\newtcolorbox[auto counter]{exo}[1][]{boxbase,colback=primary!4,colframe=primary,
  title={\textsc{Exercice}~\thetcbcounter\ifx\relax#1\relax\relax\else\textmd{ --- \itshape #1}\fi}}
\newtcolorbox[auto counter]{corrige}[1][]{boxbase,colback=excol!4,colframe=excol,
  title={\textsc{Corrigé}~\thetcbcounter\ifx\relax#1\relax\relax\else\textmd{ --- \itshape #1}\fi}}
\newcommand{\R}{\mathbb{R}}\newcommand{\N}{\mathbb{N}}\newcommand{\Z}{\mathbb{Z}}\newcommand{\ds}{\displaystyle}
% (un \usepackage{fancyhdr}+en-tête est possible ; il est ignoré côté web)
% =========================================================================
\begin{document}
\sisetup{per-mode=symbol}
\begin{exo}[dioxyde d'azote dans l'air pollué]
Une chimiste dose le dioxyde d'azote d'un air pollué : elle trouve $\qty{0,0361}{\milli\gram}$ de
\ce{NO2} dans un échantillon de $\qty{12,0}{\litre}$ d'air. Calcule la concentration $[\ce{NO2}]$ en
$\si{\gram\per\cubic\metre}$, en notation scientifique et avec le bon nombre de chiffres significatifs.
On rappelle $\qty{1}{\milli\gram}=10^{-3}\,\si{\gram}$ et $\qty{1}{\litre}=10^{-3}\,\si{\cubic\metre}$
(facteurs exacts).
\end{exo}

\begin{exo}[solubilité du chlorure de sodium]
On peut dissoudre au maximum $\qty{36}{\gram}$ de \ce{NaCl} dans $\qty{100}{\gram}$ d'eau à
$\qty{25}{\celsius}$. Quelle masse maximale de \ce{NaCl} se dissout dans $\qty{93,233}{\gram}$ d'eau à la
même température ? Donne le résultat en kilogrammes à $4$ chiffres significatifs, puis commente la
précision \emph{réellement} justifiée par les données.
\end{exo}

\begin{exo}[préparation d'une solution de \ce{NaCl}]
On dissout $m=\qty{5,00}{\gram}$ de chlorure de sodium ($M=\qty{58,44}{\gram\per\mole}$) dans l'eau, pour
obtenir $V=\qty{250,0}{\milli\litre}$ de solution.
\begin{enumerate}[label=\alph*)]
  \item Calcule la quantité de matière $n$ de \ce{NaCl}, avec le bon nombre de CS.
  \item Déduis la concentration $c=\dfrac{n}{V}$ en $\si{\mol\per\litre}$, avec le bon nombre de CS.
\end{enumerate}
\end{exo}

\begin{exo}[acidité de deux solutions]
\begin{enumerate}[label=\alph*)]
  \item Une solution d'acide chlorhydrique a $[\ce{H3O+}]=4,0\times10^{-2}\ \si{\mol\per\litre}$. Calcule
        son pH avec le bon nombre de décimales.
  \item Une autre solution a $\mathrm{pH}=5,00$. Donne $[\ce{H3O+}]$ avec le bon nombre de CS.
\end{enumerate}
\end{exo}

\end{document}
